% GENERATED FILE. Edit canonical lessons, then run npm run generate:books.
% canonical-source-hash: fnv1a64:c9295466cd163984
% canonical-lessons: FA-C512-havadar, FA-C512-varzeshkar, FA-C512-harif, FA-C512-lig, FA-C512-emtiyaz

\chapter{Teams and Fans}
\label{ch:fa-teams-and-fans}
\hlchaptermodality{512}

\begin{chapteropening}
\textbf{By the end of this chapter:} \emph{I can talk about teams, fans and points.}

Complete the last lesson of chapter 512: I can talk about teams, fans and points.
\end{chapteropening}

\section[\texorpdfstring{\fa{هوادار} (havâdâr)}{هوادار (havâdâr)}]{\fa{هوادار} (havâdâr) — a fan, a supporter}

\label{lesson:FA-C512-havadar}



\begin{quote}
Before the new one: say the Persian for walking, then the Persian for rowing.
\end{quote}

\subsection*{You'll want to know: \fa{هوادار}}
\textbf{\fa{هوادار}} (\emph{havâdâr}) — \textquotedblleft{}a fan, a supporter\textquotedblright{}.

\begin{quote}
\fa{هواداران} \fa{تیم} \fa{تا} \fa{آخر} \fa{بازی} \fa{ماندند}.
\end{quote}

\begin{quote}
\emph{havâdârân-e tim tâ âkhar-e bâzi mândand.}
\end{quote}

\begin{quote}
The team's fans stayed until the end of the match.
\end{quote}

\emph{havâ} 'desire, liking' + \emph{dâr} 'having'. Plural \emph{havâdârân}.

\begin{cousinweb}
\emph{havâdâr} is built in Persian on \emph{havâ}, 'desire, liking', a loan usually traced to Arabic, and \emph{dâr}, the stem of native \emph{dâshtan}, 'to have'. EtymDB has no entry for it.
\end{cousinweb}

\begin{grammarlens}[title={What you've built}]
One more word for telling, explaining and describing at B1.
\end{grammarlens}

\subsection*{Your turn}
\begin{itemize}
\raggedright
  \item \emph{Say it:} \emph{havâdâr}
  \item \emph{Say it:} \emph{havâdâr}, once more
  \item \emph{Say it:} a sentence of your own with \emph{havâdâr}
  \item \emph{Recall:} say the Persian for walking, then the Persian for rowing, then say \emph{havâdâr} again
\end{itemize}

\begin{tcolorbox}[breakable,colback=teal!4,colframe=teal!35!black,title={Before you move on}]
What does \fa{هوادار} mean? (\textquotedblleft{}A fan, a supporter\textquotedblright{}.) Say it once more.
\end{tcolorbox}
\section[\texorpdfstring{\fa{ورزشکار} (varzeshkâr)}{ورزشکار (varzeshkâr)}]{\fa{ورزشکار} (varzeshkâr) — an athlete, a sportsperson}

\label{lesson:FA-C512-varzeshkar}



\begin{quote}
Before the new one: say the Persian for rowing, then the Persian for a fan.
\end{quote}

\subsection*{You'll want to know: \fa{ورزشکار}}
\textbf{\fa{ورزشکار}} (\emph{varzeshkâr}) — \textquotedblleft{}an athlete, a sportsperson\textquotedblright{}.

\begin{quote}
\fa{او} \fa{ورزشکار} \fa{است} \fa{و} \fa{هر} \fa{روز} \fa{تمرین} \fa{میکند}.
\end{quote}

\begin{quote}
\emph{u varzeshkâr ast va har ruz tamrin mikonad.}
\end{quote}

\begin{quote}
She's an athlete and trains every day.
\end{quote}

\emph{varzesh} 'sport' + \emph{-kâr} 'one who does'.

\begin{cousinweb}
\emph{varzeshkâr} is built on native Persian parts: \emph{varzesh}, 'sport', from the verb \emph{varzidan}, 'to practise', and \emph{kâr}, 'work', which EtymDB traces to Middle Persian (EtymDB).
\end{cousinweb}

\begin{grammarlens}[title={What you've built}]
One more word for telling, explaining and describing at B1.
\end{grammarlens}

\subsection*{Your turn}
\begin{itemize}
\raggedright
  \item \emph{Say it:} \emph{varzeshkâr}
  \item \emph{Say it:} \emph{varzeshkâr}, once more
  \item \emph{Say it:} a sentence of your own with \emph{varzeshkâr}
  \item \emph{Recall:} say the Persian for rowing, then the Persian for a fan, then say \emph{varzeshkâr} again
\end{itemize}

\begin{tcolorbox}[breakable,colback=teal!4,colframe=teal!35!black,title={Before you move on}]
What does \fa{ورزشکار} mean? (\textquotedblleft{}An athlete, a sportsperson\textquotedblright{}.) Say it once more.
\end{tcolorbox}
\section[\texorpdfstring{\fa{حریف} (harif)}{حریف (harif)}]{\fa{حریف} (harif) — an opponent, a rival}

\label{lesson:FA-C512-harif}



\begin{quote}
Before the new one: say the Persian for a fan, then the Persian for an athlete.
\end{quote}

\subsection*{You'll want to know: \fa{حریف}}
\textbf{\fa{حریف}} (\emph{harif}) — \textquotedblleft{}an opponent, a rival\textquotedblright{}.

\begin{quote}
\fa{حریف} \fa{ما} \fa{تیم} \fa{خیلی} \fa{قوی} \fa{بود}.
\end{quote}

\begin{quote}
\emph{harif-e mâ tim-e kheyli qavi bud.}
\end{quote}

\begin{quote}
Our opponent was a very strong team.
\end{quote}

In speech: \emph{harifash nemisham} — I'm no match for him.

\begin{cousinweb}
\emph{harif} is a loan, usually traced to Arabic \emph{ḥarīf}, 'a fellow in the same trade'; Persian shifted it to 'rival'. EtymDB lists it without a source.
\end{cousinweb}

\begin{grammarlens}[title={What you've built}]
One more word for telling, explaining and describing at B1.
\end{grammarlens}

\subsection*{Your turn}
\begin{itemize}
\raggedright
  \item \emph{Say it:} \emph{harif}
  \item \emph{Say it:} \emph{harif}, once more
  \item \emph{Say it:} a sentence of your own with \emph{harif}
  \item \emph{Recall:} say the Persian for a fan, then the Persian for an athlete, then say \emph{harif} again
\end{itemize}

\begin{tcolorbox}[breakable,colback=teal!4,colframe=teal!35!black,title={Before you move on}]
What does \fa{حریف} mean? (\textquotedblleft{}An opponent, a rival\textquotedblright{}.) Say it once more.
\end{tcolorbox}
\section[\texorpdfstring{\fa{لیگ} (lig)}{لیگ (lig)}]{\fa{لیگ} (lig) — a league}

\label{lesson:FA-C512-lig}



\begin{quote}
Before the new one: say the Persian for an athlete, then the Persian for an opponent.
\end{quote}

\subsection*{You'll want to know: \fa{لیگ}}
\textbf{\fa{لیگ}} (\emph{lig}) — \textquotedblleft{}a league\textquotedblright{}.

\begin{quote}
\fa{تیم} \fa{ما} \fa{در} \fa{لیگ} \fa{اول} \fa{شد}.
\end{quote}

\begin{quote}
\emph{tim-e mâ dar lig avval shod.}
\end{quote}

\begin{quote}
Our team came first in the league.
\end{quote}

\emph{lig-e bartar} — the top league.

\begin{cousinweb}
\emph{lig} is a modern loan, usually traced to French \emph{ligue} or English \emph{league}. EtymDB has no entry for the Persian word.
\end{cousinweb}

\begin{grammarlens}[title={What you've built}]
One more word for telling, explaining and describing at B1.
\end{grammarlens}

\subsection*{Your turn}
\begin{itemize}
\raggedright
  \item \emph{Say it:} \emph{lig}
  \item \emph{Say it:} \emph{lig}, once more
  \item \emph{Say it:} a sentence of your own with \emph{lig}
  \item \emph{Recall:} say the Persian for an athlete, then the Persian for an opponent, then say \emph{lig} again
\end{itemize}

\begin{tcolorbox}[breakable,colback=teal!4,colframe=teal!35!black,title={Before you move on}]
What does \fa{لیگ} mean? (\textquotedblleft{}A league\textquotedblright{}.) Say it once more.
\end{tcolorbox}
\section[\texorpdfstring{\fa{امتیاز} (emtiyâz)}{امتیاز (emtiyâz)}]{\fa{امتیاز} (emtiyâz) — a point (in a score)}

\label{lesson:FA-C512-emtiyaz}



\begin{quote}
Before the new one: say the Persian for an opponent, then the Persian for a league.
\end{quote}

\subsection*{You'll want to know: \fa{امتیاز}}
\textbf{\fa{امتیاز}} (\emph{emtiyâz}) — \textquotedblleft{}a point (in a score)\textquotedblright{}.

\begin{quote}
\fa{تیم} \fa{ما} \fa{سه} \fa{امتیاز} \fa{گرفت}.
\end{quote}

\begin{quote}
\emph{tim-e mâ se emtiyâz gereft.}
\end{quote}

\begin{quote}
Our team got three points.
\end{quote}

Plural \emph{emtiyâzât}. It also means 'privilege, advantage'.

\begin{cousinweb}
\emph{emtiyâz} was borrowed from Arabic \emph{imtiyāz} (EtymDB), a word for distinction and privilege; Persian also uses it for a point in a game.
\end{cousinweb}

\begin{grammarlens}[title={What you've built}]
One more word for telling, explaining and describing at B1.
\end{grammarlens}

\subsection*{Your turn}
\begin{itemize}
\raggedright
  \item \emph{Say it:} \emph{emtiyâz}
  \item \emph{Say it:} \emph{emtiyâz}, once more
  \item \emph{Say it:} a sentence of your own with \emph{emtiyâz}
  \item \emph{Recall:} say the Persian for an opponent, then the Persian for a league, then say \emph{emtiyâz} again
\end{itemize}

\begin{tcolorbox}[breakable,colback=teal!4,colframe=teal!35!black,title={Before you move on}]
What does \fa{امتیاز} mean? (\textquotedblleft{}A point (in a score)\textquotedblright{}.) Say it once more.
\end{tcolorbox}
